\section{Changing the averaging parameters without paying for depth}\label{cs:section}

The finite reservoir construction provides nearly optimal trials with one
shared sampling level. Later intermediate trials can retain shared sampling
variables at several levels. The next two sections describe and then remove
that sharing. Before introducing their
structural bookkeeping, we need an estimate that controls a simultaneous
change of the averaging parameters. Estimating one level at a time would
charge for the number of levels, which is not known when the approximation
tolerances are chosen. The estimate below avoids that dependence.

Fix the root data of a finite-depth trial and work under its conditional
law. Let
$\mathcal F_0\subseteq\cdots\subseteq\mathcal F_d$ be its history
filtration, with $\mathcal F_0$ trivial under this conditional law. For a
bounded $\mathcal F_d$-measurable variable $X$ and parameters
$0<t_1\le\cdots\le t_d\le1$, define
\[
 Y_d=X,\qquad
 Y_{l-1}=\frac1{t_l}\log\E[e^{t_lY_l}\mid\mathcal F_{l-1}],
 \quad 1\le l\le d,
 \qquad \mathcal T_{\mathbf t}X=Y_0.
\]
The conditional sampling kernels are fixed while the parameters vary.
For a single unconditional operation, write
$T_tV=t^{-1}\log\E e^{tV}$ for $t>0$ and $T_0V=\E V$.

\begin{lemma}[Sensitivity independent of depth]\label{cs:sensitivity}
Suppose $|X|\le C$ almost surely, where $C\ge0$. For any two
nondecreasing vectors $\mathbf t,\mathbf t'\in(0,1]^d$,
\begin{equation}\label{cs:exponent-bound}
 |\mathcal T_{\mathbf t}X-\mathcal T_{\mathbf t'}X|
 \le K(C)\|\mathbf t-\mathbf t'\|_\infty,
 \qquad K(C)=\frac52 C^2e^{4C}.
\end{equation}
For a single operation and $0\le t\le1$,
\begin{equation}\label{cs:small-exponent}
 0\le T_tX-\E X\le\frac12 C^2e^{2C}t.
\end{equation}
Both assertions hold conditionally on almost every root value. Their
constants are independent of the depth and of the fixed sampling kernels.
\end{lemma}

\begin{proof}
The case $C=0$ is immediate. In general every $Y_l$ lies in $[-C,C]$.
Conditional Jensen gives nonnegative drifts
\[
 D_l=Y_{l-1}-\E[Y_l\mid\mathcal F_{l-1}].
\]
Their expectations telescope to
$\sum_l\E D_l=Y_0-\E X\le2C$. Since a conditional mean minimizes mean
square distance to a constant,
\begin{align*}
 \Var(Y_l\mid\mathcal F_{l-1})
 &\le\E[(Y_l-Y_{l-1})^2\mid\mathcal F_{l-1}]\\
 &=\E[Y_l^2\mid\mathcal F_{l-1}]-Y_{l-1}^2+2Y_{l-1}D_l\\
 &\le\E[Y_l^2\mid\mathcal F_{l-1}]-Y_{l-1}^2+2CD_l.
\end{align*}
Taking expectations and summing gives the depth-independent variance
budget
\begin{equation}\label{cs:variance-budget}
 \sum_{l=1}^d\E\Var(Y_l\mid\mathcal F_{l-1})
 \le\E X^2-Y_0^2+2C(Y_0-\E X)\le5C^2.
\end{equation}
This expectation is under the original conditional root law.

We now calculate the effect of one parameter. For a fixed bounded random
variable $V$, let $\psi(t)=\log\E e^{tV}$. Under exponential tilting by
$e^{tV}/\E e^{tV}$, one has $\psi''(t)=\Var_t(V)$. Hence
\[
 \frac{\dd}{\dd t}T_tV
 =\frac{t\psi'(t)-\psi(t)}{t^2}
 =\frac1{t^2}\int_0^t u\Var_u(V)\dd u.
\]
If $|V|\le C$ and $0\le u\le1$, the tilted density is at most $e^{2C}$.
Using the untilted mean as a competitor for the tilted variance gives
$\Var_u(V)\le e^{2C}\Var(V)$. Therefore
\begin{equation}\label{cs:local-derivative}
 0\le\frac{\dd}{\dd t}T_tV\le\frac12e^{2C}\Var(V).
\end{equation}
The same estimate holds conditionally at each history.

When $t_l$ varies, $Y_l$ is fixed because it uses only the deeper
parameters. Each outer operation averages the resulting change against
its normalized exponential density. The chain rule therefore gives
\[
 \frac{\partial Y_0}{\partial t_l}=\E[W_{l-1}g_l],\qquad
 g_l=\left.\frac{\dd}{\dd t}
 \left(\frac1t\log\E[e^{tY_l}\mid\mathcal F_{l-1}]\right)
 \right|_{t=t_l},
\]
where $W_0=1$ and
\[
 W_b=\exp\left(\sum_{j=1}^b t_j(Y_j-Y_{j-1})\right).
\]
These are densities relative to the original history law; the effective
values $Y_j$ need not be independent. Summation by parts gives
\begin{align*}
 \log W_b
 &=t_bY_b-t_1Y_0-\sum_{j=1}^{b-1}(t_{j+1}-t_j)Y_j\\
 &\le C\left(t_b+t_1+\sum_{j=1}^{b-1}(t_{j+1}-t_j)\right)
 =2Ct_b\le2C.
\end{align*}
The nondecreasing order of the parameters is used in this inequality.
Combining it with Equations~\eqref{cs:local-derivative}
and~\eqref{cs:variance-budget} yields
\[
 \sum_{l=1}^d\left|\frac{\partial Y_0}{\partial t_l}\right|
 \le\frac12e^{4C}\sum_{l=1}^d
       \E\Var(Y_l\mid\mathcal F_{l-1})
 \le\frac52C^2e^{4C}.
\]
The straight segment joining $\mathbf t$ and $\mathbf t'$ consists of
nondecreasing vectors in $(0,1]^d$. Integrating the directional
derivative proves Equation~\eqref{cs:exponent-bound}. Equal adjacent
parameters cause no problem: the nested transforms are differentiable
in every positive coordinate, and the density estimate holds along the
entire segment.

For one operation, Equation~\eqref{cs:local-derivative} and
$\Var(X)\le C^2$ give the upper bound in
Equation~\eqref{cs:small-exponent} after integration from $0$ to $t$.
Jensen gives the lower bound, and bounded convergence gives
$T_tX\to\E X$ as $t\downarrow0$. Boundedness justifies all conditional
differentiations. The argument holds under each conditional root law for
which the kernels are defined, hence for almost every root value.
\end{proof}

Two consequences will be used after the structural part. First, consider
pairwise disjoint consecutive blocks
\[
 I_j=\{b_j+1,\ldots,c_j\}\subseteq\{1,\ldots,d\},
 \qquad t_{c_j}-t_{b_j+1}\le\eta.
\]
Replace all parameters in $I_j$ by $t_{c_j}$. The new vector is still
nondecreasing and differs by at most $\eta$ in supremum norm, so the
change of $\mathcal T X$ is at most $K(C)\eta$, regardless of the number
of blocks. In a block with common parameter $s$, the tower property
combines its transforms into a single conditional operation:
\[
 e^{sY_{b_j}}=\E[e^{sY_{c_j}}\mid\mathcal F_{b_j}].
\]

Second, fix initial data and suppose that a random variable $z$ and a
random vector $u$ are independent under the resulting conditional law.
For $V(z,u)\in[-C,C]$ and $0<m<m_0\le1$, independence gives
\[
 T_m^{(z,u)}V=T_m^z(T_m^uV),
 \qquad
 0\le T_m^{(z,u)}V-\E_z[T_m^uV]
 \le\frac12C^2e^{2C}m_0.
\]
Here the superscript specifies the variables integrated. The inner
quantity $T_m^uV$ lies in $[-C,C]$, so the inequality is precisely
Equation~\eqref{cs:small-exponent} applied to the outer operation. It
replaces one small-parameter logarithmic average over $z$ by ordinary
expectation while retaining the logarithmic average over $u$.

For either endpoint of a packet with $D$ clauses, the leaf loss has
absolute value at most $\log2+\beta D$ by
Equation~\eqref{ctr:packet-bound}. With $D\le D_0$, both consequences
therefore apply with $C=\log2+\beta D_0$. They compare one packet at
each endpoint. The removal argument will reduce a long horizon to one
packet before invoking them.
